\cotefile{72}

\scanpage{98}

\uline{Dém.}

(1°) Pour la formule (5), qui se réduit à

\[s_i\in\unsure{E_i|\mathcal E_i}\qquad 0\leq i\leq n+1\]

on procède par réc. sur $i$, la formule étant triviale pour $i=0$. On prouve en $=$ temps (6), (qui est \unsure{vide|vides|vraie} pour $i=0$ si on pose $a_{-1}=s_0$, $\lambda_{-1}=1$, $E_{-1}=0$). On a

\[\begin{aligned}s_{i+1}&=\tau_i s_i=\underset{\substack{\rotatebox{-90}{$\in$}\\ E_i}}{s_i}+\langle s_i,a'_i\rangle a_i\in E_i+k a_i=E_{i+1},\\s_{i+1}&\equiv\langle s_i,a'_i\rangle a_i\quad(E_i)\end{aligned}\]

Comme

\[s_i\equiv\deleted{\unsure{\ell|k}}\lambda_{i-1}a_{i-1}\quad(E_{i-1})\]

on aura

\[\langle s_i,a'_i\rangle=\lambda_{i-1}\langle a_{i-1},a'_i\rangle=\lambda_i,\quad\deleted{\unsure{E}}\]

d'où la formule $(*)$ pour $i+1$ au lieu de $i$, si on suppose $a'_i\mid E_{i-1}=0$, OK.

\uline{Corollaire} Soient

\[\left\{\begin{aligned}h_i&=k a_i\\ H_i&=\operatorname{Ker}a'_i\end{aligned}\right.\quad\bigg|\quad\deleted{\unsure{0}}-1\leq i\leq n\deleted{\unsure{,\ h_{-1}=ks}.}\]

\margin{\uline{NB} $h_i,H_i$ \deleted{\unsure{ont|sont}} \unsure{définis|défini} en termes de $\tau_i$, \uline{dès que} $\{\unsure{a_i|s_i}\}$ est libre}

les sous-variétés associés à $\tau_i$, posons aussi $\left\{\begin{aligned}a_{-1}&=s_0\\h_{-1}&=ks_0\\[-2pt]&\deleted{\unsure{a_{-1}=s_0}}\end{aligned}\right.$. Les conditions \deleted{équivalentes} ci-dessous (A$'$) et (B$'$) sont équivalentes et impliquent la formule (15) $(E_i=k(s_0,\ldots,s_{i-1}))$

(A$'$) On a

\[\deleted{\unsure{(A')}}\quad h_i\subset H_j\quad\text{si}\quad -1\leq i<j\leq n,\quad j\ne i+1\tag*{(15$'$)}\]

\[\begin{aligned}&\deleted{\unsure{E_i=k(s_0,\ldots,s_{i})\quad 0\leq i\leq n+1}}\\ & (s_i)_{0\leq i\leq n+1}\text{ est libre }\bigl(\uline{\text{ou encore}},\ \deleted{\unsure{\text{équivalent}}}\\ &\deleted{\unsure{\text{il suffit de les}}}\ \deleted{\unsure{\Leftrightarrow H_j,\ h_j,\ i=n+1}}\ \deleted{\unsure{i,j}}\ h_i\not\subset \unsure{H'_{i+1}|H_{i+1}}\text{ en chaque}\\ &\text{point }-1\leq i\leq n-1\quad\text{indépendants}\bigr)\end{aligned}\tag*{(16$'$)}\]

(B$'$) On a

\[\langle a_i,a'_j\rangle=0\quad\text{pour}\quad-1\leq i<j\leq n,\quad j\ne i+1\tag*{(17)}\]

\[\deleted{\lambda_n\in k^*}\quad(a_i)_{-1\leq i\leq n}\text{ est libre}\tag*{(18)}\]

\[\lambda_n\in k^*\tag*{(19)}\]

\margin{\deleted{\unsure{$a_i$}}\quad\unsure{$0$|$a$}}

% margin note written sideways

En effet, (17) n'est qu'une reformulation de (15$'$), et il reste à prouver, en supposant (15$'$) et (17) vérifiées, que \unsure{(16$'$)|(15$'$)} équivaut à (18) + (19). \deleted{Pour $n=1$ \unsure{on c'est évident}.} On peut supposer que \deleted{\unsure{l'on a}} (si $n\geq1$) \inserted{$\lambda_{\unsure{n-1|n-2}}\ne0$}, c'est

% archive sheet number: 97
\scanpage{99}

déjà prouvé pour $n-1$, et que $s_0\ldots s_n$ et $(a_{\unsure{-1|-9}},\ldots,a_{n-1})$ sont lin. indép. et $(s_0,\ldots s_n)$ engendrent $\unsure{E_n|\mathcal E_n}$. \deleted{\unsure{Or on a|On a donc} \unsure{par la}}

Dire $(s_0,\ldots s_n,s_{n+1})$ libre \unsure{ssi} $s_{n+1}$ libre mod $E_n$ \deleted{$s_{n+1}\equiv\lambda_n a_n\bmod E_n$, donc $s_{n+1}$ \unsure{est|et} $(s_0\ldots s_{n+1})$ libre \unsure{ssi}} \deleted{$s_{n+1}$} or par \unsure{1°)|(6)} on a

\[s_{n+1}\equiv\lambda_n a_n\quad(E_n)\]

donc la condition équivaut à $\alpha)$ $a_n$ libre mod $E_n$, i.e. (18), et $\beta)$ \deleted{\unsure{$\lambda_n$}} la condition (19). \deleted{Pour}

Donc on a l'équivalence dite, et de plus ces \deleted{$E_n=k(s_0,\ldots,s_n$\unsure{$\ E_n$}} conditions équiv. impliquent

\[E_{n+1}=E_n+ks_{n+1}=k(s_0,\ldots,s_n)+ks_{n+1}=k(s_0,\ldots s_{n+1}),\]

ce qui achève la dém. du corollaire.

\textcircled{\scriptsize $2^\circ$} On procède par récurrence sur $n$, supposons (\unsure{si|ici} $n\ne0$) que c'est prouvé déjà pour $n-1$. Donc on peut supposer que ttes les conditions écrites dans A, B sont satisfaites pour le système $(\deleted{\unsure{E}}s_0,\tau_0,\ldots\tau_{n-1})$, \inserted{et qu'on a \deleted{\unsure{donc} $E_i=k(s_0\ldots s_{i-1}$} de plus (14) pour $0\leq i\leq n$,} et prouver qu'alors le syst. de conditions

\[(\mathrm A'')\quad\left\{\begin{aligned}\tau_n&\text{ commute à }\tau_0\ldots\tau_{n-2}&&\alpha).\\\tau_n s_0&=s_0 &&\beta)\\s_{n+1}&\text{ libre mod }E_n\ \deleted{\unsure{\}}}&&\gamma)\end{aligned}\right.\]

équivaut au syst. de conditions

\[(\mathrm B'')\quad\left\{\begin{aligned}\left\{\begin{aligned}\langle a_i,a'_n\rangle&=0\\ {}a'_i,a_n\rangle&=0\end{aligned}\right\}\quad\text{pour }0\leq i\leq n-2&&\big\}\ \alpha)\\ \langle s_0,a'_n\rangle=0&&\beta)\\ a_n\text{ libre mod }E_n&&\gamma)\\ \lambda_n\in k^*&&\delta)\end{aligned}\right.\]

et que ces conditions impliquent

\[E_{n+1}=k(s_0,\ldots s_{n+1}).\tag*{(20)}\]

% archive sheet number: 98
\scanpage{100}

Comme $s_{n+1}=\tau_n s_n=s_n+\langle s_n,a'_n\rangle a_n$, \inserted{la condition de $(\mathrm A'')$:} $s_{n+1}$ libre mod $\unsure{E_n|\mathcal E_n}$, implique $a_n$ libre mod $E_n$, donc $(s_0,a_0,\ldots,a_n)$ libre.

\deleted{En particulier, \unsure{si} $(a_i,a_n)$ est libre si $0\leq i\leq n-1$, donc $\tau_i$ commute à $\tau_n$ ssi $\langle a'_i,a_n\rangle=\langle a_i,a'_n\rangle=0$. De même $(s_0,a_n)$ étant libre, la condition $\tau_n s_0=s_0$ \unsure{La première et troisième condition de} \unsure{impliquent|implique}}

Donc on peut supposer, pour établir l'équivalence de $(\mathrm A'')$ et $(\mathrm B'')$, que $(s_0,a_0,\ldots,a_n)$ est libre. Mais alors \deleted{\unsure{tous}} $(a_i,a_n)$ étant libre pour $0\leq i\leq n-2$, la condition \deleted{$\alpha)$ de} $(\mathrm A''\alpha)$ s'interprète par $(\mathrm B''\alpha)$, $(\mathrm A''\beta)$ par $(\mathrm B''\beta)$, il reste à montrer que $(\mathrm A''\gamma)$ équivaut à $(\mathrm B''\delta)$, \inserted{\uline{moyennant} les \deleted{\unsure{deux}} conditions précédentes supposées satisfaites.} Or $(\mathrm A'',\gamma)$, par l'hyp. faite de liberté de $a_n$ mod $E_n$, équivaut à

\[\langle s_n,a'_n\rangle\in k^*,\]

or par hyp. de récurrence et 1°) déjà prouvé, on a

\[s_n\equiv\lambda_{n-1}a_{n-1}\quad\bmod E_{\unsure{n-1|n=1}}\]

et $a'_n\mid E_{n-1}=0$ par $(\mathrm B''\alpha,\beta)$, donc

\[\langle s_n,a'_n\rangle=\lambda_{n-1}\langle a_{n-1},a'_n\rangle=\lambda_n\]

donc on trouve bien la condition $\lambda_n\in k^*$ de $(\mathrm B''\delta)$. Enfin il est évident que ttes ces conditions impliquent $E_n+ks_{n+1}=E_n+ka_n\unsure{=|\equiv}E_{n+1}$, et comme $E_n=k(s_{\unsure{0|1}},\ldots s_n)$ par hyp. de récurrence, on en déduit bien (20),

cqfd.

% archive sheet number: 99
\scanpage{101}

\textcircled{\scriptsize $3^\circ$} On a déjà prouvé (14), il faut prouver (15) (16).

\uline{Formule} (15) : Pour $j<i$ cela provient de $\tau_j=\mathrm{id}+a'_j\otimes a_j$, et $a_j\in\unsure{E_i|\mathcal E_i}$ \inserted{(vrai sous les hyp. A) B))}, pour $j>i$ cela provient de $a'_j\mid E_i=0$, qui implique $E_i\subset\underset{\substack{\Vert\\(E^{\tau_j})}}{H_j}$.

\uline{Formule} 16\qquad On a

\[E_i=E_{i-1}+ka_{i-1}\]

\[\begin{aligned}\text{donc}\qquad\tau_i E_i&=\underset{\Vert}{\uline{\tau_i E_{i-1}}}+k\tau_i a_{i-1}\\ &=E_{i-1}+k\underset{\substack{\rotatebox{-90}{$\in$}\\ k^*}}{\underline{\langle a_{i-1},a'_i\rangle}}a_i=E_{i-1}+k\deleted{\unsure{E}}a_i\end{aligned}\]

et comme $(a_{i-1},\deleted{\unsure{et}}a_i)$ \deleted{\unsure{sont}} est libre mod $E_{i-1}$, cela prouve notre assertion.

\uline{Corollaire 2} Sous les conditions du \deleted{\unsure{th.}} \inserted{corollaire} \inserted{(a fortiori sous les conditions du th 1 2°)}) \deleted{\unsure{on aura|on a}}

\[\deleted{\unsure{H_i\overset{\mathrm{dfn}}=|H_i\overset{\mathrm{df}}=}}\quad\left\{\begin{aligned}E^{\tau_i}&=\operatorname{Ker}a'_i\ (\overset{\mathrm{dfn}}=H_i)\\ka_i&=\operatorname{Im}(\tau_i-\mathrm{id})\ (\overset{\mathrm{dfn}}=h_i)\end{aligned}\right.\quad\text{pour}\quad0\leq i\leq n\tag*{(21)}\]

\[E_i\subset H_j\quad\deleted{\unsure{(\ =E^{\tau_j})\ \operatorname{Ker}a'_i|(\operatorname{Ker}a'_j)}}\qquad\text{pour }0\leq i<j\leq n\deleted{\unsure{a}}\tag*{(22)}\]

\[E_i\cap H_i=E_{i-1},\tag*{(23)}\]

\[\tau_i\tau_{i+1}\ne\tau_{i+1}\tau_i,\quad\uline{\text{même modulo }k^*}\qquad\deleted{\unsure{\text{si}}}\ (\text{si }k\ne0)\tag*{(24)}\]

\[\deleted{\underset{\substack{\Vert\\ka_i}}{h_i}\subset\underset{\substack{\Vert\\\operatorname{Ker}a'_{i+1}}}{H_{i+1}}}\quad\text{i.e.}\quad\langle a_i,a'_{i+1}\rangle\in k^*\ \inserted{-1\leq i\leq n}\ (\text{donc }\ne0\text{ si }k\ne0)\tag*{(25)}\]

\[\underset{\substack{\Vert\\ka_i}}{h_i}\not\subset\underset{\substack{\Vert\\\operatorname{Ker}a'_{i+1}}}{H_{i+1}}\qquad\text{si}\quad k\ne0\tag*{(25$'$)}\]

La formule \unsure{(22)|(21)} est déjà prouvée, (21) 1° provient du fait que $\{a_i\}$ est libre, (21 2°) du fait que $a'_i:\deleted{\unsure{\text{est surj}}} E\longrightarrow k$ est surjectif, \deleted{\unsure{ceci}} car $\langle a_{i-1},a'_i\rangle\in k^*$ (NB $a_{-1}=s_0$ !) grâce à (13) (4). \deleted{On a} La formule (23) provient de $E_i=k(a_{-1},\ldots,a_{i-1})$, et que $H_i=\operatorname{Ker}a'_i$, $a'_i$ \unsure{nul|nulle} sur $a_j$ $(-1\leq j\leq i-2)$, $\langle a_{i-1},a'_i\rangle\in k^*$.

% archive sheet number: 100
\scanpage{102}

Pour (24), on note que

\[\tau_i\tau_{i+1}-\underset{\rho\in k^*}{\rho}\tau_{i+1}\tau_i=\underset{\substack{\rotatebox{-90}{$\in$}\\ k^*}}{\underline{\langle a_{i+1},a'_i\rangle}}\,a'_{i+1}\otimes a_i-\rho\underset{\substack{\rotatebox{-90}{$\in$}\\ k^*}}{\langle a_i,a'_{i+1}\rangle}\,a'_i\otimes a_{i+1}\]

\deleted{qui \unsure{reste|doit} est $\ne0$ car \unsure{d'après les hyp.} la famille de $\{a'_{i+1}\otimes a_i,a'_i\otimes a_{i+1}\}$ est \uline{libre} dans $\operatorname{End}(\unsure{E|\mathcal E})$ --- son image par $\operatorname{End}(E)\longrightarrow E$, $u\longmapsto u(x)$ est formée de $\{a'_{i+1}(x)a_i,a'_i(x)a_{i+1})$, qui \unsure{son} est libre si}

Son image est \unsure{formé|formée} des

\[\langle a_{i+1},a'_i\rangle a'_{i+1}(\unsure{n|x})a_i-\rho\langle a_i,a'_{i+1}\rangle a'_i(x)a_{i+1},\]

je dis que c'est $k(a_i,\unsure{a_{i+1}|h_{i+1}})=h_i+h_{i+1}$, \deleted{ce} qui \deleted{\unsure{provient de|prouve}} signifie aussi que

\[(a'_i,a'_{i+1}):E\unsure{\text{---}|\longrightarrow} k^2=ke_1\oplus ke_2\]

est surjectif, en effet \deleted{l'image de} \inserted{sa restriction} $k(a_{i-1},\deleted{\unsure{a|\oplus}}a_i)$ a la matrice \deleted{$a_{i-1}$ est $(\underset{\substack{\rotatebox{-90}{$\in$}\\ k^*}}{\langle a_{i-1},a'_i\rangle},0)$, donc l'image de $E$ contient $ke_1$, enfin l'image}

\[\begin{pmatrix}\langle a_{i-1},a'_i\rangle&*\\0&\langle a_i,a'_{i+1}\rangle\end{pmatrix}\]

de déterminant $\langle a_{i-1},a'_i\rangle\langle a_i,a'_{i+1}\rangle\in k^*$, en vertu de (25), lui-$=$ \unsure{traduction|conséquence} évidente de (13) $\lambda_n\in k^*$, compte tenu de (4). Enfin (25$'$) est une conséquence immédiate de (25),

cqfd.

% archive sheet number: 101
\scanpage{103}

\uline{Remarques} Si $\unsure{E|\mathcal E}$ est loc. libre de rg $n+2$, alors on trouve des conditions équivalentes $(\overline{\mathrm A})$, \deleted{\unsure{$\overline{\mathrm B}$|$(\mathrm B)$}} et $(\overline{\mathrm B})$, en renforçant (A), (B), remplaçant \deleted{\unsure{(8)}}(9) de (A) et (12) de (B) respectivement par

\[(s_i)_{0\leq i\leq n+1}\text{ \unsure{forme|forment} une base de }E\tag*{(9$'$)}\]

\[(a_i)_{-1\leq i\leq n}\text{ \unsure{forme|forment} une base de }E.\tag*{(12$'$)}\]

\deleted{Quand ces conditions} C'est aussi équivalent à la conjonction de (A), ou (B), avec la condition

\[E_{n+1}=E\tag*{(26)}\]

\deleted{Soit alors $a'_{0}=a'_{-1}$} \deleted{Si ces conditions sont satisfaites, \unsure{complétons} \unsure{Soit} Soit \unsure{$a'_j$|$a'_{-1}$} \inserted{$s'_0=a'_{-1}$}, une forme lin. sur $E$ telle que}

\[\deleted{\operatorname{Ker}a'_{-1}=k(a_0,\ldots,a_n)}\tag*{\deleted{(27)}}\]

\deleted{elle est donc déterminée, modulo scalaire $\in k^*$, par la seule donnée des $\tau_i$ $(0\leq i\leq n)$, du}

Idem pour cor. 1 [conditions $(\overline{\mathrm A}'),(\overline{\mathrm B}')$, équivalent à (A$'$) ou (B$'$) $+{}$(26)]

\uline{Proposition} \unsure{Soit|Soient} $E$ un $k$-module loc. libre de rg $n+2$, $\tau_i$ $(0\leq i\leq n)$ des pseudo-réflexions dans $E$,

\[\tau_i=\mathrm{id}+a'_i\otimes a_i\qquad 0\leq\unsure{i|0}\leq n,\tag*{(27)}\]

\deleted{satisfaisant (28) $\tau_i\tau_j=\tau_j\tau_i$ dans $\underline{\mathrm P}(E)$.}

Pour qu'il existe \inserted{localement} \unsure{un|des} $s_0\in E$ tel que les conditions $(\overline{\mathrm A})(\overline{\mathrm B})$ du Th. 1 \inserted{(remarque ci-dessus)} soient satisfaites (conditions \inserted{commutation et} de ``non dég''), il f et s.

% archive sheet number: 102
\scanpage{104}

que les conditions suivantes soient satisfaites

a) $(a_i)_{0\leq i\leq n}$ est \deleted{libre} strictement libre dans $\unsure{E|\mathcal E}$ (i.e. engendre un sous-Module facteur direct de rg $n+1$)

a$'$) $(a'_i)$ est strictement libre dans $E^\vee$.

b) $\langle a_i,a'_j\rangle=0\qquad 0\leq i<j\leq n,\ j\ne i+1$\quad (ou encore $h_i\subset H_j$ pour ces \unsure{$(i,j)$|$i,j$})

c) \deleted{\unsure{$H_j\ H_i$}} $\langle a_i,a'_{i+1}\rangle\in k^*\qquad(0\leq i\leq n-1)$

\[(\Longleftrightarrow h_i\not\subset H_{i+1}\text{ en chaque fibre},\quad0\leq i\leq n-1)\]

(Ces conditions impliquent \deleted{donc} (28) $\langle a_i,a'_j\rangle=0$ pour $0\leq i,j\leq n$, $i,j$ non consécutifs). Pour qu'on ait

\uline{Dém} La nécessité de a) b) \inserted{c)} est claire, \deleted{\unsure{pour la|prouvons}} \inserted{prouvons} celle de a$'$), \deleted{\unsure{pour cela il faut|prouvons} il suffit de prouver une propriété de stabilité des conditions du Th. 1 par dualité. Or considérons $E_n\subset E_{n+1}=E$, qui est un hyperplan, et soit $s'_0$ une}

On suppose donc qu'on a $s_0$ satisfaisant les conditions dites, on pose. \deleted{$a_{-1}=s_0$, et on définit}

\[a_{-1}=s_0\in E,\qquad a'_{n+1}\in E^\vee\]

défini par

\[\operatorname{Ker}a'_{n+1}=k(a_{-1},a_0\ldots,a_{n-1})\]

\[\text{donc}\quad\left\{\begin{aligned}\langle a_i,a'_{n+1}\rangle&=0&&-1\leq i\leq n-1\\ \inserted{c_n\overset{\mathrm{dfn}}{=}}\langle a_n,a'_{n+1}\rangle&\in k^*\end{aligned}\right.\]

\[\begin{aligned}\text{Considérons}\quad &(a_{-1},a_0,\ldots a_n)\in E^{n+2}\\ &(a'_0,a'_1,\ldots,a'_{n+1})\in E^{\unsure{n+2|-+2|\vee\,n+2}}\end{aligned}\]

et la matrice qu'ils forment, c'est une matrice triangulaire, avec sur la diagonale

% archive sheet number: 103
\scanpage{105}

\[c_{-1}=\langle a_{-1},a'_0\rangle,\quad c_0=\langle a_0,a'_1\rangle\ldots,\ \deleted{\unsure{c}}c_n=\langle a_n,a'_{n+1}\rangle\]

donc de déterminant égal à

\[\deleted{\unsure{a|\text{et}}}\ \Delta=c_{-1}c_0\ldots c_n=\underset{\substack{\rotatebox{-90}{$\in$}\\ k^*}}{\lambda_n}\underset{\substack{\rotatebox{-90}{$\in$}\\ k^*}}{c_n}\ \in k^*,\]

ce qui prouve que \deleted{\unsure{$\langle a_0$|les}} $(a'_0\ldots a'_n\ a'_{n+1})$ forme une base de $\unsure{E^\vee|\mathcal E^\vee}$, d'où (a$'$).

Inversement, supposons a) a$'$) b) c) \unsure{satisfait|satisfaites|satisfaits}, cherchons $s_0\in E$ satisfaisant, avec $\tau_0,\ldots\tau_n$, aux conditions de $(\overline{\mathrm B}')$. \inserted{La condition $(\overline{\mathrm B}',17)$ pour $i\ne-1$ n'est autre que b).}

\deleted{\unsure{La Tout d'abord|En Tout d'abord} (28) s'interprète, par le calcul habituel, par \unsure{ce qui précède|Rappelons} \unsure{la condition|la situation} (17) $\langle a_i,a'_j\rangle=0\quad 0\leq i,j\leq n$, $i,j$ non conséc.}

\deleted{La condition (7) est déjà satisfaite par hypothèse, la condition (8) signifie $\langle s_0,a'_i\rangle a_i=0$ pour $1\leq i\leq n$, et comme $a_i$ est libre, cela équivaut à}

\[\deleted{\langle s_0,a'_i\rangle=0\qquad1\leq i\leq n}\tag*{\deleted{(29)}}\]

-- En vertu de a$'$), l'ens des $s_0$ satisfaisant à \deleted{ces conditions} \inserted{(17) pour $i=-1$ (\unsure{posant|puisque|pour} $s_0=a_{-1}$)} est un sous-module facteur direct, loc. libre de rang $(n+2)-n=2$; \deleted{\unsure{soit}} \deleted{Soit} soit $F$. \inserted{Pour tenir compte de $(\overline{\mathrm B}',19)$ pour $i=0$,} Soit $D\subset F$ la droite, définie par les conditions $\langle s_0,a'_i\rangle=0\quad0\leq i\leq n$ dans $E$, \deleted{donc} ou la seule condition $\langle s_0,a'_0\rangle\deleted{\unsure{\in}}=0$ \deleted{\unsure{$\in$|$k^*$}} dans $F$. La condition (19) $\lambda_n\in k^*$ signifie compte tenu de \unsure{(c)|(6)} et (4), que $\langle s_0,a'_0\rangle\in k^*$, \deleted{\unsure{$\langle s_0,a'_0\rangle C$|$\langle s_0,a'_0\rangle\in$}} i.e. $s_0\notin D$ sur chaque fibre.

% archive sheet number: 104
\scanpage{106}

Reste à exprimer \inserted{$(\overline{\mathrm B}')$} (18), pour ceci considérons \deleted{\unsure{pour|la} condition \unsure{forme lin.|formes lin.} $a'_{-1}$ qui s'annule sur les $a_i$ \unsure{l'hom.|le module}} $L=\unsure{E|\mathcal E}/k(a_0,\ldots,a_n)$, et l'hom. induit par $E\unsure{\text{---}|\longrightarrow}L$, $F\longrightarrow L$, \deleted{cet hom} je dis que c'est un épi, i.e. que c'est un épi sur \inserted{chaque} fibre \deleted{\unsure{de $f$|point}}. On est ramené au cas d'un corps de base, et à prouver que dans ce cas on n'a pas $F\subset k(a_0,\ldots,a_n)$, \deleted{ou encore que \unsure{la fo|Or}}

Or $F\cap k(a_0,\ldots,a_{\unsure{n+1|n-1}})=(0)$, car le système de $n$ formes linéaires $a'_1,\ldots,a'_n$ sur $k(a_{\unsure{0|j}}\ldots a_n)$ a une matrice triangulaire, de déterminant le produit des $c_i=\langle a_i,a'_{i+1}\rangle$ $(0\leq i\leq n-1)$, donc $\in k^*$. Comme $k(a_0,\ldots,a_{\unsure{n|n-1}})$ est de codim 1 dans $k(a_0,\ldots,a_n)$, \deleted{\unsure{cela ex}} et $\dim F=2$, cela implique qu'on ne peut avoir $F\subset k(a_0,\ldots a_n)$.

Soit \deleted{$D'$} $\Delta=\operatorname{Ker}(F\unsure{\longrightarrow|\text{---}}L)$, qui est une deuxième droite de $F$. La condition \inserted{$(\overline{18})$} sur $s_0$ signifie que $s_0\notin\Delta$ en chaque fibre. Donc la conjonction de $(\overline{17})$ $(\overline{18})$ $(\overline{19})$ signifie que $s_0$ est une section de $F$ \inserted{(fibré vect. de rg 2)} qui n'est dans $D$ ni $\Delta$ en aucune fibre -- il en existe tjrs \unsure{loc.\,lmt|localement}, OK.

% archive sheet number: 105
\scanpage{107}

Classes de Stiefel-Whitney et classe de Brauer de la forme invariante

% archive sheet number: 106
\scanpage{108}

\uline{\unsure{Classes|Classe} de Stiefel-Whitney}.

Supposons $f_V$ lisse \deleted{et} 2 inv, i.e. \deleted{2}

\[\delta(f_V)=\deleted{\unsure{2}}\bigl(\beta_0^n\beta_1^{n-1}\ldots\beta_{n-1}\bigr)\rho\quad\text{inv. et 2 inv.}\]

Calculons $w_*(f_V)$. Pour ceci, on va procéder par récurrence, en calculant l'orthogonal de $\unsure{\underline O_S|\mathcal O_S}[a_1,\ldots,a_n]$, \textsuperscript{*} comme $\psi_V(a_i)=-(\beta_0\ldots\beta_{i-1})a'_i$, c'est aussi l'orthogonal de $a'_1\ldots a'_n$, donc \inserted{engendré par} la projection \deleted{de} \inserted{$s_0-\frac1\rho\pi$} $s_0=a_{-1}$ sur $V$, dans $\unsure{E|\mathcal E}=V\oplus\underline O_S\pi$, soit $p=\mathrm{id}-\frac1\rho\pi'\otimes\pi$ la dite projection. [notée $\widetilde s_0$]

\margin{$^*$ Mais il faut supposer $f\mid\underline O_S(a_1,\ldots a_n)$ lisse, donc $\sigma=\rho_{\unsure{1n|1,n|n}}$ inversible}

\[\left[\begin{aligned}&\widetilde s_0=-\frac1\rho\bigl(\unsure{\zeta_0|\xi_0}a_0+\zeta_1a_1\ldots+\zeta_na_n\bigr)\\ &\left\{\begin{aligned}\zeta_n&=1\\\zeta_{n-1}&=2\\ &\text{---}\\\zeta_i&=L_{n-i}(\beta_{i+1},\ldots,\beta_{n-1})=\rho_{i+1,n-1}\\ &\text{---}\\\zeta_0&=L_n(\beta_1\ldots\beta_n)=\rho_{1,n}=\sigma\end{aligned}\right.\end{aligned}\right]\]

\[\begin{aligned}f_V(\widetilde s_0)&=f_E(s_0-\tfrac1\rho\pi)=\underbrace{f_E(s_0)}_{0}+\frac1{\rho^2}\underbrace{f_E(\pi)}_{-\widetilde\rho\widetilde\sigma}\deleted{\unsure{+}}-\frac1\rho\underbrace{\varphi(s_0,\pi)}_{\underbrace{\langle s_0,\underbrace{\psi(\pi)}_{-\sigma\pi'}\rangle}_{-\sigma}}\\ &=-\frac{\widetilde\rho\widetilde\sigma}{\rho^2}+\frac{\sigma}{\rho}=\frac1{\rho^2}\bigl[\underbrace{\rho\sigma}_{2\widetilde\rho\widetilde\sigma}-\widetilde\rho\widetilde\sigma\bigr]=\frac{\widetilde\rho\widetilde\sigma}{\rho^2}\end{aligned}\]

\[\partial\bigl(f_V(\widetilde s_0)\bigr)=\partial(\widetilde\rho\widetilde\sigma)=\partial\bigl(\widetilde\rho_{0,n-1}\widetilde\rho_{1,n-1}\bigr)\]

\margin{Bockstein $H^0({},\mathbb G_m)\longrightarrow H^1(-,\unsure{\mu_m|\mu_n|\mu_2})$}

% archive sheet number: 107.
\scanpage{109}

Donc on trouve

\[w_*(f_{V_{n+1}}(\beta_0,\ldots,\beta_{n-1}))=\bigl[1+\underset{\substack{\rotatebox{-90}{$\in$}\\ H^1(-,\deleted{\unsure{\mu_2|\mu}}\mathbb Z/2)}}{\underline{\partial(\widetilde\rho_{0,n-1}\widetilde\rho_{1,n-1})}}\bigr]\,w_*(f_{V_n}(\beta_{\unsure{1|{}}},\ldots,\beta_{n-1}))\]

On trouve ainsi de proche en proche, pourvu que les calculs aient un sens i.e. que \deleted{tous} les $\rho_{i,n-1}$ soient inversibles $(0\leq i\leq n-1)$

\[\begin{aligned}w_*(f_{V_{n+1}}(\beta_0,\ldots,\beta_{n-1}))={}&[1+\partial(\widetilde\rho_{0,n-1}\widetilde\rho_{1,n-1})][1+\partial(\widetilde\rho_{1,n-1}\widetilde\rho_{2,n-1})]\\ &\ldots[1+\partial(\widetilde\rho_{\deleted{\unsure{n-2,n-1}}\inserted{n-2,n-1}}\widetilde\rho_{\deleted{\unsure{n-2,n-1}}\inserted{n-1,n-1}})]\,\unsure{\bar w_*|w_*}(f_{V_2}(\beta_{n-1}))\end{aligned}\]

De $=$, on trouve,

\[\begin{aligned}w_*(f_{V_2}(\beta_{n-1}))&=\deleted{\bigl(1+\partial[\underbrace{\widetilde\rho_{n-1,n-1}}_{4-\beta_{n-1}}\underbrace{\widetilde\rho_{n,n-1}}_{\Vert\,1}]\bigr)\underbrace{\bigl(1+\partial(\underbrace{\widetilde\rho_{n,n-1}}_{1}\underbrace{\widetilde\rho_{n+1,n}}_{1})\bigr)}_{\unsure{1}}}\\ &\bigl(1+\partial(\beta(4-\beta))\bigr)=[1+\partial(2+\alpha)\deleted{\unsure{+}}(2-\alpha))]\end{aligned}\]

En effet, faisant le calcul pour $f_V(\beta)$ sur $a_0,a_1$ on trouve, si on veut

\[\left\{\begin{aligned}f(a_0)&=1\qquad\partial f(a_0)=0\\a_0^\perp&=\{x\unsure{,|{}}a_0+ya_1\mid 2x-\beta y=0\quad\text{i.e. }\deleted{y=}x=\tfrac\beta2y\}\end{aligned}\right.\]

engendré par $\beta a_0+2a_1=u$

\[\begin{aligned}f(u)&=\deleted{\unsure{\beta^2+4\beta-|\beta^2+4\beta}}\ \beta^2\underbrace{f(a_0)}_1+4\underbrace{f(a_1)}_\beta+2\beta\underbrace{\varphi(a_0,a_1)}_{-\beta}\\ &=4\beta-\beta^2=\beta(4-\beta)\end{aligned}\]

d'où \deleted{$w_*($} la formule annoncée.

% archive sheet number: 108
\scanpage{110}

On trouve, en résumé

\[w_*(f_V(\beta_0,\ldots,\beta_{n-1}))=\left(\prod_{0\leq i\leq n-2}(1+\partial(\widetilde\rho_{i,n-1}\widetilde\rho_{i+1,n-1}))\right)\left(1+\partial\bigl(\underbrace{\beta_{n-1}(4-\beta_{n-1})}_{(2+\alpha_{n-1})(2-\alpha_{n-1})}\bigr)\right)\]

formule valable si tous les $\rho_{i,n-1}$ $(0\leq i\leq n-1)$ sont inversibles, ainsi que les $\beta_i$ $(0\leq i\leq n-1)$

Bien sûr, on aurait pu conduire le calcul en sens inverse, en mettant à part $(a_0\ldots a_{n-1})$ etc, \deleted{on} trouve-t-on d'autres expressions remarquables ?

Calculons $w_*(f_{\unsure{E|\mathcal E}})$, \deleted{si on suppose} en notant que $V^\perp=\unsure{\underline O_S|\mathcal O_S}\pi$ et $f(\pi)=-\widetilde\rho\widetilde\sigma$ et on trouve donc un facteur supplémentaire $1+\partial(\underbrace{-\widetilde\rho\widetilde\sigma}_{-\widetilde\rho_{0,n-1}\widetilde\rho_{1,n-1}})$. Posons donc

\[\begin{aligned}\partial(-1)&=\varepsilon&&&&\text{éléments de}\\ \partial(\widetilde\rho_{i,n-1})&=\unsure{\eta_i|y_i}&&0\leq i\leq n-1&& H^1(-,\mathbb Z/2)\\\partial(\widetilde\beta_{n-1})&=\eta_n\end{aligned}\]

On trouve

\[w_*(f_V)\deleted{\unsure{)}}=\prod_{0\leq i\leq n-1}(1+\eta_i+\eta_{i+1})\qquad\text{NB } w_{n+1}(f_V)=0\ !\]

\[w_*(f_E)=(1+\varepsilon+\eta_0+\eta_1)\prod_{0\leq i\leq n-1}(1+\eta_i+\eta_{i+1})\]

% archive sheet number: 109
\scanpage{111}

Si $n=2$, ceci donne

\[\left\{\begin{aligned}\varepsilon&=\partial(-1)\\ \unsure{\eta_0|y_0}&=\partial(\widetilde\rho)=\partial(\underbrace{4-\beta_0-\beta_1}_{-(\alpha_0+\alpha_1)})=\partial[-(\alpha_0+\alpha_1)]\\ \eta_1&=\partial\widetilde\sigma=\partial(4-\beta_1)=\partial(2-\alpha_1)\\ \eta_2&=\partial\beta_1=\partial(2+\alpha_1)\end{aligned}\right.\]

\[w_* f_V=(1+\eta_0+\eta_1)(1+\eta_1+\eta_2)\]

avec hyp. parasite que $\widetilde\sigma\deleted{\unsure{\ne}}=2-\alpha_1$ \deleted{\unsure{et $\beta_1=2+\alpha_1$}} \inserted{soit} inversible. Un calcul plus satisfaisant, sans cette hypothèse, s'obtient en prenant l'orthogonal de $\unsure{\underline O_S|\mathcal O_S}(a_0,a_2)$ engendré par $u=\beta_0 a_0+2a_1+a_2$, on trouve $f(u)=-\beta_0(\alpha_0+\alpha_1)$, et $f(a_0)=1$, $f(a_2)=\beta_0\beta_1$ syst. \unsure{orthonormal|orthonormé|orthogonal}. On pose alors

\[\begin{aligned}\varepsilon&=\partial(-1)\\\eta_0&=\partial(\beta_0)\\\eta_1&=\partial(\beta_1)\\\eta_2&=\partial[-(\alpha_0+\alpha_1)]\end{aligned}\]

et on trouve

\[\begin{aligned}w_*(f_V)&=(1+\eta_0+\eta_2)(1+\eta_0\deleted{\unsure{)}}+\eta_1)\qquad\text{i.e.}\\w_1(f_V)&=\deleted{\unsure{\eta_0+}}\eta_1+\eta_2\\w_2(f_V)&=\deleted{\unsure{\eta_1^2}}(\eta_0+\eta_1)(\eta_0+\eta_2)=\eta_0^2+(\eta_0\eta_1+\eta_1\eta_2+\eta_2\eta_0)\end{aligned}\]

% archive sheet number: 110
\scanpage{112}

On a

\[c(f_V)\overset{\mathrm{dfn}}=w_2(f_V)+w_1(f_V)^2=(\unsure{\eta_0|y_0}+\eta_1+\eta_2)^2+(\eta_0\eta_1+\eta_1\eta_2+\eta_2\eta_0)\]

(Son annulation est néc et suff pour que le fibré en droites projectives envisagé soit isomorphe à celui défini \deleted{par \unsure{l'équat.|l'équation}} \inserted{par la} forme $x^2+y^2+z^2$). On a

\[\operatorname{Br}(f_V)\overset{\mathrm{dfn}}=c(f_V)-\underbrace{c(\deleted{\unsure{f_{2E}}}\,X^2+yz)}_{\varepsilon^2}=(\varepsilon+\eta_0+\eta_1+\eta_2)^2+(\eta_0\eta_1+\eta_1\eta_2+\eta_2\eta_0)\]

L'annulation de cette expression est néc et suff pour que le fibré en droites projectives \deleted{\unsure{envis.|associé}} défini par $f_V$ soit associable à un fibré vect. de rang 2, en fait, c'est la classe de Brauer de ce fibré. -- Le calcul universel se fait sur

\[S_0=\operatorname{Spec}\mathbb Z[\beta_0,\beta_1]_{2\beta_0\beta_1(4-\beta_0-\beta_1)}\]

On \unsure{notera|note} que dans $H^1(S_0,\mathbb Z/2)$, les \deleted{sections} \inserted{éléments} $\varepsilon,\eta_0,\eta_1,\eta_2$ forment une base. Il est plausible que les 10 expressions $\varepsilon^2$, $\eta_0^2$, $\eta_1^2$, $\eta_2^2$, $\varepsilon \eta_0$, $\varepsilon \eta_1$, $\varepsilon \eta_2$, $\eta_0\eta_1$, $\eta_1\eta_2$, $\eta_2\eta_0$ sont

% archive sheet number: 111
\scanpage{113}

\margin{?}

lin. indép. dans $H^2(S_0,\mathbb Z/2)$.\ldots

\uline{Pb} Considérons une forme quadratique \inserted{lisse} à coeff. indéterminés, à coefficients dans $\mathbb Z[(\unsure{a_{ij}|c_{ij}})_{1\leq i\leq j\leq n}]_{2\delta^{\prime}((c_{ij}))}$

``Calculer'' sa classe de Stiefel-Whitney totale ?

\uline{NB} Le $H^1(S_0,\mathbb Z/2)$ \inserted{$\simeq H^0(S_0,G_m)/2$} est engendré par $\varepsilon=\partial(-1),\ \deleted{\unsure{\partial|\xi}}\underbrace{\partial(2)}_{\xi},\ \underbrace{\partial(\delta'(c_{ij}))}_{\unsure{\eta|y}=w_1(f)}$, dont les 2 premiers deviennent nuls sur une extension biquadratique $\deleted{\unsure{\mathbb Q|\mathrm Q|\mathbb Z}}\mathbb Z_2[\sqrt{-1},\sqrt{2}]$ -- a fortiori sur $\mathbb C$, le deuxième donne lieu à $w_1$. Il est \inserted{donc} clair que multiplicativement à partir de ces éléments, on ne peut exprimer \inserted{aucun des} $w_2,w_3$ \ldots{} etc.

% archive sheet number: 112
